\section{The Question Answered}

\subsection{Separating the claim into its parts}

The sentence under test bundles four assertions, and they do not stand or fall
together. Separating them is most of the work.

\begin{studyframe}\small
\textbf{(i)} That the burial cloths were left in a particular and noticeable
arrangement.

\textbf{(ii)} That Peter was the one who noticed it.

\textbf{(iii)} That the arrangement was evidence from which something could be
inferred.

\textbf{(iv)} That Peter drew the inference, and that what he inferred was the
Resurrection.
\end{studyframe}

\textbf{(i) is what the text says}, in every Greek manuscript tradition, in the
Latin the Church read for fifteen centuries, and in the English translated from
that Latin. The linen cloths were lying; the face-cloth was not lying with
them, but apart, in a wrapped-up state, in one place. That is not an inference;
it is the verse.

\textbf{(ii) is also what the text says}, and this is the part of the popular
claim that is usually defended badly and happens to be right. Verses 6 and 7
form one sentence whose subject is Simon Peter and whose verb is
\emph{the\=orei}. The face-cloth and everything said about it lie inside what
Peter beholds. The other disciple, in verse 5, saw the \emph{othonia} from
outside and did not go in; the evangelist says so twice. On the evidence of the
Fourth Gospel, the man who noticed how the burial cloths were left was Peter.
Chrysostom saw this and made a point of it---the evangelist ``witnesses to the
exactness of Peter's search''---and Theophylact and Aquinas repeat him. It is a
real feature of the text and the claim is entitled to it.

\textbf{(iii) is defensible, ancient, and not certain.} Chrysostom's case that
the arrangement tells against theft is a real argument with three independent
steps---motive, physical difficulty, exposure---and it is directed against an
explanation the Gospels themselves record as current. Theophylact reproduces it;
Aquinas endorses it by name. But it was made by a minority of the commentators
who had the occasion, it depends on a reconstruction of how the spices were
used, and it eliminates theft rather than removal. The strongest form of it is
not Chrysostom's but the modern one---that the cloths had collapsed where the
body had been and the face-cloth still held the shape of a head---and that form
is an inference from what the text does not say as much as from what it does.

\textbf{(iv) is what neither evangelist says, and what one of them denies.}

\subsection{Project synthesis}

\begin{studybox}{Project synthesis: does the premise survive?}
\textbf{Judgment.} The maintainer's premise survives in part and fails in its
operative half. Peter is, on the Fourth Gospel's own showing, the disciple to
whom the observation of the cloths is attributed: he entered first, he saw
most, and the description of the face-cloth is grammatically inside what he
beheld. But no evangelist attributes the inference to him. John assigns the
believing of Jn 20:8 to the other disciple, in the singular, by name-in-effect;
Luke has Peter leave the tomb \emph{thaumaz\=on}, wondering, a participle Luke
pairs twenty-nine verses later with \emph{apistount\=on}, not believing. The
insight, so far as Scripture assigns one, belongs to the beloved disciple. What
belongs to Peter is the looking.

\textbf{Reasons.} (1) Jn 20:8 has a singular verb and a named subject that is
not Peter, in a sentence the evangelist could have written in the plural and did
not. (2) Jn 20:9, joined by \emph{gar}, says of both men that they did not yet
know the scripture that he must rise---which is why Augustine, Gregory, and
Theophylact all lower the content of the believing rather than extend its
subject. (3) Luke, independently or in the earliest surviving compression of
John, ends Peter's visit in wonder, not faith. (4) The Johannine pattern of
delayed understanding (Jn 2:22; 12:16) is the evangelist's own explanation of
how his disciples receive events. (5) Every harmonization that would deliver
the claim either concedes it (belief came later), or costs more than it buys
(deleting Lk 24:12 supplies a better witness against the claim than it
removes), or requires reading one Father's construal of John against another
evangelist's plain sense.

\textbf{The strongest counterargument, at full strength.} It is not a
sceptical argument but a patristic one, and it is considerable. Two Greek
Fathers of the first rank read verse 8 collectively and evidentially.
Chrysostom: ``From this they believed in the Resurrection.'' Cyril of
Alexandria, more explicitly still: the two disciples, being competent witnesses
because ``two in number, as the Law enjoined,'' did not yet meet the risen
Christ ``but infer His Resurrection from the bundle of linen clothes, and
henceforth believed that He had burst asunder the bonds of death.'' Ambrose, in
Latin, says flatly that Peter could not have doubted. Chrysostom generalizes his
reading into an account of Easter faith itself: the apostles ``from the proof of
the napkins \ldots\ straightway received the word concerning the Resurrection,
and before they saw the body, exhibited all faith.'' And Aquinas, having given
Augustine's reading as the resolution, sets Chrysostom's beside it---\emph{vel,
secundum Chrysostomum} \ldots\ \emph{credidit, scilicet vera fide}---with a
reading of the syntax of verse 9 that would carry it, and does not close it off.
On this account the maintainer's sentence is not a popular error at all: it is
the Antiochene and Alexandrian reading of the passage, and the man who
transmitted it to the Latin West was Thomas Aquinas.

\textbf{Why the counterargument does not carry, and what it does establish.} It
does not carry, for three reasons. Chrysostom's and Cyril's collective reading
is reached by expounding John continuously without Luke in view; neither
addresses Peter's wondering, because neither is commenting on Luke at that
point. Both pass Jn 20:9 without the weight Augustine puts on it---Cyril in
fact reverses it, having the disciples read the events ``in the light of the
prophecies,'' which is close to the opposite of what the verse says. And even
at its strongest the counterargument delivers (iv) for \emph{both} disciples on
the strength of the \emph{othonia}; Cyril never mentions the face-cloth, and the
peculiar arrangement of the \emph{soudarion}---the whole point of the popular
sentence---plays no part in his argument at all.

What the counterargument does establish is that the maintainer's premise is a
respectable minority reading of the Fourth Gospel with excellent patristic
credentials, not a modern devotional invention. That distinction should be kept
sharp, because the other claim examined in this study is a modern devotional
invention, and the two should not be filed together.

\textbf{What would change this judgment.} A textual witness of weight reading
the verbs of Jn 20:8 in the plural would settle it the other way at a stroke;
none is known to this study, and the apparatus of the registered Greek edition
records no variant there. Failing that, a demonstration that Lk 24:12 is
neither Lucan nor an early compression of John---that is, a late and dependent
gloss with no independent standing---would remove the only text that says
anything at all about Peter's state of mind, and would leave the field to the
Chrysostom--Cyril reading of a Gospel that never denies Peter's faith, only
declines to assert it. Conversely, the judgment would be strengthened past
argument by any early witness that read Jn 20:8 as this study reads it
\emph{and} engaged Luke while doing so; no such witness was located.
\end{studybox}

\subsection{Why the attribution to Peter arose}

If the insight belongs to the beloved disciple, the popular sentence needs
explaining, and four causes are visible in the evidence gathered here. They are
offered as this study's reconstruction.

\textbf{The text gives Peter the observation.} This is the largest cause and the
most respectable. A reader of Jn 20:3--10 comes away with a vivid picture of
Peter inside the tomb looking at things, and a very thin picture of the other
disciple, who is described almost entirely by what he did not do---did not go
in, did not know the scripture. The verse that assigns the believing is seven
words long and has no object. It is easy to read the paragraph and remember
Peter.

\textbf{Luke 24:12, whatever its origin, concentrates the episode on Peter.}
Westcott and Hort's account of the verse, if right, makes it the earliest
surviving instance of the move: a reader condensing John 20:3--10 ``with
omission of all that relates to `the other disciple'.'' Whether that is what
happened or whether Luke wrote it, the canonical result is the same---one Gospel
in which the tomb visit is Peter's alone.

\textbf{Chrysostom's collective reading became the tradition's other voice, and
Aquinas carried it.} A reading in which ``they believed'' is available in the
\emph{Catena aurea} and in the \emph{Lectura super Ioannem}, in a commentary
that has just spent a paragraph on Peter's diligence, is one short step from a
reading in which Peter believed. The step is short enough to be taken without
noticing.

\textbf{And Peter is the one people remember.} A story about the first Pope
noticing something at the empty tomb has a shape that a story about an unnamed
disciple does not. This is not a criticism of anyone; it is how narrative
memory works, and it is worth saying that the tradition's own instinct here is
not simply wrong. Peter did see. He is the reason we know about the face-cloth
at all.

\subsection{What is left, and it is not small}

Stripping the sentence of what it cannot carry does not leave nothing. It leaves
this.

An eyewitness tradition, written down late in the first century, reports that
two men went into a tomb on a Sunday morning and found the wrappings of a corpse
without the corpse; that the head-cloth was not with the rest; that it was
separate, and in a place of its own, and---by the choice of a verb its own
author elsewhere uses for swathing a body---still in the condition of something
that had been wound round a head.

Whatever else is true, somebody thought that arrangement worth remembering for
sixty years and worth a whole verse when he finally wrote it down. He does not
say what he made of it. He says only that the other man went in, and saw, and
believed. The Gospel's reticence about what exactly was seen and exactly what
was believed has kept the tradition arguing for sixteen centuries---Chrysostom
against Augustine, Cyril against Gregory, Aquinas declining to choose---and that
argument is a better thing to inherit than a manufactured custom about napkins.

The verse is stranger than the legend. It has the advantage of being there.
